% !TEX root = ../algebra_02.tex

\section{Конечные подгруппы мультипликативной группы поля}


\begin{Theorem}{О конечной подгруппе мультипликативной группы поля}{}
	Пусть $K$ --- поле, а $G < K^*$ --- конечная подгруппа его мультипликативной группы (без нуля). Тогда $G$ является циклической группой.
\end{Theorem}

\begin{Remark}{Связь с кольцом вычетов}{}
	Из Основного курса алгебры (ОМА) известно, что мультипликативная группа кольца вычетов $(\mathbb{Z}/n\mathbb{Z})^*$ является циклической тогда и только тогда, когда $n = p^k$, $n = 2p^k$ (где $p$ --- нечетное простое, $k \in \mathbb{N}$), либо $n = 1, 2, 4$.
\end{Remark}

\begin{proof}
	Так как $G$ --- конечная абелева группа, по структурной теореме она изоморфна прямому произведению примарных циклических групп:
	\[
	G \cong \prod_{i=1}^l \mathbb{Z}/p_i^{m_i}\mathbb{Z}.
	\]
	Предположим от противного, что группа $G$ не является циклической. Это означает, что в её разложении среди простых оснований $p_i$ есть одинаковые. Без ограничения общности (WLOG), пусть $p_1 = p_2 = p$.
	
	Перейдем временно к аддитивной записи группы $G'$ (где $G' \cong G$).
	В первой компоненте разложения $\mathbb{Z}/p^{m_1}\mathbb{Z}$ рассмотрим элементы вида $\overline{k \cdot p^{m_1-1}}$, где $0 \le k \le p-1$. Умножение любого такого элемента на $p$ дает $\overline{0}$. Всего таких элементов $p$ штук.
	
	Аналогично, во второй компоненте $\mathbb{Z}/p^{m_2}\mathbb{Z}$ элементы вида $\overline{l \cdot p^{m_2-1}}$, где $0 \le l \le p-1$, при умножении на $p$ дают $\overline{0}$. Их также $p$ штук.
	
	Значит, в прямом произведении элементы вида:
	\[
	(\overline{k \cdot p^{m_1-1}}, \overline{l \cdot p^{m_2-1}}, \overline{0}, \dots, \overline{0})
	\]
	будут давать $\overline{0}$ при умножении на $p$. Общее количество таких элементов составляет не менее $p \times p = p^2$:
	\[
	\left| \{ g' \in G' \mid p \cdot g' = 0 \} \right| \ge p^2.
	\]
	
	Возвращаясь к мультипликативной записи в исходной подгруппе $G < K^*$, мы получаем, что множество элементов $x$, порядок которых делит $p$ (то есть $x^p = 1$), состоит как минимум из $p^2$ элементов:
	\[
	\left| \{ x \in G \mid x^p = 1 \} \right| \ge p^2 \implies \left| \{ x \in K^* \mid x^p = 1 \} \right| \ge p^2.
	\]
	
	Это означает, что многочлен $x^p - 1 = 0$ имеет в поле $K$ не менее $p^2$ корней. Однако мы знаем, что над любым полем многочлен степени $p$ может иметь не более $p$ корней.
	Мы пришли к противоречию ($p^2 \le p$), следовательно, наше предположение неверно. Все простые основания $p_i$ в прямом произведении различны.
	
	Раз все $p_i$ различны, то порядки сомножителей попарно взаимно просты, а значит, прямое произведение изоморфно одной большой циклической группе:
	\[
	G \cong \mathbb{Z} \left/ \left( \prod_{i=1}^l p_i^{m_i} \right) \mathbb{Z} \right..
	\]
	Таким образом, $G$ --- циклическая группа. Теорема доказана.
\end{proof}